problem:  
Let \(M^n\) be a closed, simply connected Riemannian manifold with **nonnegative sectional curvature**, and suppose \(G\) is a compact Lie group acting **isometrically and fixed point homogeneously** on \(M\). Let \(F \subset \mathrm{Fix}(G, M)\) be a fixed point component of maximal dimension and let \(H\) be the principal isotropy subgroup.  

Known results imply that \(M\) admits a **double disk-bundle decomposition** \(M = D(F) \cup_\phi D(N)\), where \(N\) is a submanifold corresponding to the opposite boundary stratum in the orbit space \(M/G\). The gluing map \(\phi\colon \partial D(F)\to \partial D(N)\) encodes the **equivariant topology** of \(M\).  

**Problem:**  
Assume that the **representation of \(H\)** on the normal sphere to \(F\) is **non-transitive** on the unit sphere, and that \(\pi_1(G/H)\) is **nontrivial**.  
Can one determine whether the diffeomorphism class of \(M\) is uniquely determined by the isotropy data \((G,H,F)\) and the homotopy class of the gluing map \(\phi\)?  
Equivalently, is there rigidity of the gluing map under nonnegative curvature—i.e., does nonnegative curvature constrain \(\phi\) uniquely up to isotopy?  

Why is it a "good" problem:  
This problem isolates an **as-yet unexplored rigidity phenomenon** at the intersection of **nonnegative curvature geometry**, **transformation group theory**, and **topological classification**.  

- **Conceptual core:** Known structural theorems (Grove–Searle, Grove–Wilking) guarantee the existence of a double disk-bundle decomposition for fixed point homogeneous manifolds, but they do not describe the possible **gluing maps** \(\phi\). These gluing maps encode global topological complexity. Controlling or classifying them under curvature assumptions would extend fixed point homogeneity from a coarse topological statement to a new **geometric rigidity principle**.  

- **Innovative focus:** Introducing a non-transitive isotropy representation and nontrivial \(\pi_1(G/H)\) highlights settings beyond previously classified, highly symmetric examples (e.g., \(S^n\), projective spaces). Here the problem asks whether curvature restricts the topological “twist” in forming \(M\) from its bundle pieces — an unexplored invariant possibly linking Alexandrov geometry and smooth topology.  

- **Simplicity and depth:** The question is defined by a small number of geometric inputs (\(G, F, H, \phi\)) yet invokes deep analysis of how **curvature, group action geometry, and bundle attachment data interact**. It is simple to state but probes an intricate frontier of geometric topology.  

- **Cross-field significance:** Solutions might require blending techniques from **equivariant Ricci flow**, **Alexandrov geometry**, and **surgery theory**—a step toward understanding when curvature enforces topological rigidity in the presence of symmetry.  

Thus this problem is “good” because it asks a sharp, unstudied, and conceptually profound question: *does nonnegative curvature rigidify the topological gluing freedom in fixed point homogeneous manifolds with nontransitive isotropy?*